\section{Realigning and propagating the marked comparison}
\label{sec:allarity-marked-propagation}

The finite polynomial comparison must survive each block integration.
At an intermediate level the marked expansion contains the desired
principal insertion, a small regular remainder, and a forced remainder.
We explain how its leading pieces are reassigned to the prescribed
supports and how the resulting errors accumulate. The identities use
the actual representation at that level.

\subsection{Why a connected principal component stays in one cell}

Let $\rho\ge0$ be an elementary slice radius. Each elementary principal
piece has a nonempty set of marked reader sites; each such reader belongs
to its full support, and the full support is within distance $\rho$ of
every one of those readers. Assume distinct
replica cells have reader separation strictly greater than $2\rho$, and
\[
                  (3+2n)\rho\le dB.
\]
Consider a connected component in the graph joining elementary pieces
whose full supports intersect. Two marked readers on the same piece
are at distance at most $\rho$, since a reader lies in that piece's
support. Readers on two adjacent pieces have distance at most
$2\rho$, by using a point in the intersection. Both bounds are at most
$2\rho$, so the readers belong to the same replica cell. Following a
path of intersecting pieces shows that all marked readers of the
component belong to one cell.

There are at most $n$ distinct marker labels in that cell. Make a graph
on the labels used by the component, joining labels appearing on the
same or on adjacent pieces. The preceding support path gives a walk
between any two of these labels. Erase cycles from the walk. The
remaining simple path uses at most $n-1$ edges, each of geometric
length at most $2\rho$. Thus the distance between any two marked
readers is at most $2n\rho$. Every point in the component's full
support is at distance at most $\rho$ from one of its readers, so
it is at distance at most $(2n+1)\rho\le dB$ from every marked
reader of the component. Its marker diameter is also at most $dB$.
This proves the required enlarged-support bounds independently of
how many elementary pieces the component contains.

For the selected construction the elementary slice radius is $\rho=3B$.
The purchase $6n+3\le a/(3+2n)$ gives the final distance bound
$(2n+1)\rho=(6n+3)B\le aB/(3+2n)$; it does not construct
the slices. In particular it supplies the bound $aB$ required for
the enlarged full supports in the regrouping below. The selected
replica separation is stronger than $2\rho$. Subsequent regrouping of pure principal
components preserves their full supports and marker unions, so the
same argument applies to the principal connected logarithm.

\subsection{Marker prices and spatial prices}

Fix $n\ge1$ marks in each replica cell. Every principal piece has a
nonempty marker subset $I\subseteq\{1,\ldots,n\}$ and a prescribed
spatial support $X_{c,I}$ in cell $c$. The geometric hypotheses are the
following: pieces with the same marker mode have a unique prescribed
support; each actual principal support is contained in that support;
and two prescribed supports containing the same spatial site belong
to the same replica cell. These statements are made for full inspected
supports, including the masks. Let
\[
 \operatorname{chg}_e(X)=\exp(e\operatorname{cov}_B(X)),\qquad
 \operatorname{chg}_e(X_{c,I})\le R_e,
\]
where $e\ge0$ is the support exponent, $B$ the current radius and
$\operatorname{cov}_B$ its covering number. Use one constant marker
weight $\omega>0$ at every mark, of every cell, replica and period.
The weight of a mode $I$ is $\omega^{|I|}$. We do not assume
$\omega\le1$; in the application below $\omega\ge1$. Since
$I\subseteq\{1,\ldots,n\}$, every mode satisfies
\begin{equation}\label{eq:allarity-mode-weight}
 \omega^{|I|}\le\bar\omega_n,\qquad
 \bar\omega_n=\max(1,\omega)^n .
\end{equation}
For $\omega\le1$ this is the bound $\omega^{|I|}\le1$.

If each prescribed piece has weighted coefficient price at most $Z$,
its total incidence price at a fixed mark $(c,j)$ is at most $2^nZ$.
Indeed only subsets $I$ of that same cell can contain $(c,j)$, and
there are at most $2^n$ subsets. This estimate is independent of the
number of replica cells.

A second estimate converts incidence price to spatial price. Suppose
actual pieces $s_a$ have nonnegative prices $p_a$ and nonempty marker
supports, and the total price of pieces containing any fixed mark is
at most $N$. Assign piece $a$ to its prescribed support
$X_{c(a),I(a)}$. At a spatial site $v$, if any such support contains
$v$, all such supports lie in one cell $c$. Charge each summand to
one of its nonempty marks and then overcount by summing over every
mark in that cell. It follows that
\begin{equation}\label{eq:allarity-spatial-incidence}
 \sum_{a:\,v\in X_{c(a),I(a)}}
            p_a\operatorname{chg}_e(X_{c(a),I(a)})
                    \le n N R_e.
\end{equation}
This is why the spatial charge is proportional to $n$, while the
number of principal subset types is bounded by $2^n$.

\subsection{Replacing the leading pieces on their prescribed supports}

Write the current regular source as a sum of remainder pieces $H_a$
and leading pieces $S_b$. Let $F_{c,I}$ be the desired principal source
at mode $(c,I)$, with support $X_{c,I}$. Group the $S_b$ by this marker
mode and define on its healthy region
\[
 L_{c,I}(U)=\sum_{b:\,\operatorname{mode}(S_b)=(c,I)}S_b(U)
                        -F_{c,I}(U).
\]
Assume $|L_{c,I}(U)|\le y$ there. Attach the full support $X_{c,I}$
to $L_{c,I}$ and apply its healthy mask. The incidence price of these
new regular remainders is at most
\begin{equation}\label{eq:allarity-realign-regular}
                         2^n \bar\omega_n R_e y.
\end{equation}
Indeed the new piece at mode $(c,I)$ has value at most $y$ on its
healthy region, support charge at most $R_e$ and marker weight
$\omega^{|I|}$, so its weighted price is at most
$\omega^{|I|}R_ey\le\bar\omega_nR_ey$ by
\eqref{eq:allarity-mode-weight}. The preceding subset count applies.

The mask change also has a forced part. A leading piece may be healthy
on its smaller support while its prescribed larger support is not
healthy. Split its original masked value into its value on the larger
healthy region and the complementary part. The latter is supported
where the enlarged support contains a flag. Insert these complementary
pieces into the forced alphabet, keeping their original marker modes.
The identity is an identity of functions on every field: on the larger
healthy region it is the defining equality for $L_{c,I}$; on its
complement the difference has just been assigned to the forced part.
Apply the exact regular/forced regrouping of
Section~\ref{sec:soft-log-fusion} to the enlarged supports. No division
by an activity value is involved.

Here is its quantitative form. Suppose the previous forced spatial
price is $s$, the leading incidence price is $Z$, and the new forced
inflation exponent is $g\ge0$. If
\begin{equation}\label{eq:allarity-realign-forced-input}
 e^g s+n Z R_{e+1}\le\delta,
           \qquad (2B+1)^4\delta\le1,
\end{equation}
then the regrouped forced price is at most
\begin{equation}\label{eq:allarity-realign-forced-output}
                           e^{-g}(e^\delta-1).
\end{equation}
Here is the proof for the marker-valued factors. Work in the
commutative algebra generated by the marker variables $z_j$ with
$z_j^2=0$. For $I$ a subset of markers write $z_I=\prod_{j\in I}z_j$.
The norm $\|\sum_I a_I z_I\|_\omega=\sum_I |a_I|\omega^{|I|}$ is
submultiplicative for every $\omega>0$: a pair of monomials with
overlapping marker sets multiplies to zero, and otherwise, for
disjoint $I,J$, the weight $\omega^{|I\cup J|}=\omega^{|I|}\omega^{|J|}$
is the product of their weights. No bound $\omega\le1$ is used. Every moved piece is a single $\ell_j=z_{I_j}h_j$ with
$I_j\ne\varnothing$. Thus $\ell_j^2=0$ and $\exp(\ell_j)=1+\ell_j$.
The mask difference defining $h_j$ vanishes whenever the enlarged
support is healthy, so a nonzero moved factor has a flag in its full
support.

Expand the exponential of the sum of the moved pieces as the product
of the factors $1+\ell_j$ and multiply by the old forced gas. In a
nonzero summand each moved factor meets an old forced support: its flag
is covered by the old forced inventories. Taking connected components
of the full-support intersection graph and flattening them back
therefore gives the bijection of
Section~\ref{sec:soft-log-fusion}. This argument preserves the
inventory and the product in the marker algebra. Overlapping marker
sets contribute zero on both sides. The expansion is justified by
absolute convergence in the weighted norm, since the constituent
prices are summable on the finite spatial lattice.

Assign inflated weights $e^g p_i$ to old forced pieces of price $p_i$
and unchanged weights to the moved pieces. At exponent $e+1$ the
site sum is at most $\Gamma=e^g s+nZR_{e+1}$ by
\eqref{eq:allarity-spatial-incidence}. Every surviving connected
component contains an old forced piece, so its product norm is
bounded by $e^{-g}$ times the product of the inflated weights.
The distinct connected-set bound \eqref{eq:ope-forest-price},
with reserve one and counting factor $(2B+1)^4$, bounds the new
site price at exponent $e$ by $e^{-g}(e^\Gamma-1)$.
Since $\Gamma\le\delta$, this proves
\eqref{eq:allarity-realign-forced-output}. The regrouping does not
change the archived background: setting all markers to zero kills
every moved factor, leaving the original background at its original
address.
Thus a regular remainder price $z$ becomes at most
$z+2^n\bar\omega_nR_e y$, by \eqref{eq:allarity-realign-regular},
while the forced price changes as displayed, and the polynomial
density is unchanged.

\subsection{The recurrence along the comparison interval}

Fix a top comparison depth $k$, an intermediate level $l$, a larger
arity allowance $u\ge\max(2,n)$, and a monitor $H>1$ fixed by $k$.
Here $k,l,u,n$ are natural numbers with $l\le k$, and all scalar
price bounds are nonnegative.
Let the tolerance exponent be $q\ge6$ and the regular step multiplier
be $\alpha\ge0$. Set
\[
 d_0=H^{-q},\qquad d_j=(\alpha+2)^j d_0.
\]
The principal sources are the connected polynomials of the same
Gaussian history at the current level, with the same cell and marker
labels. At the terminal comparison level their higher connected
coefficients vanish: the formal logarithm of
$\prod_i(1+z_iP_i)$ is $\sum_i z_iP_i$. Its singleton terms are
exactly the launched calibrated insertions. Consequently the initial
marked representation has regular remainder price $d_0$, after the
specified terminal-source purchase, which must also bound its initial
forced price by $G_k$.

Suppose the actual one-step marked compiler at each level has the
following quantitative outputs when its incoming remainder price is
$z\ge d_0$. Before realignment it supplies remainder price at most
$(\alpha+1)z$, principal incidence price at most $H^{-2}$, and forced
price at most $G_i^2$ at the coarse level. Here
$G_i=e^{-g_i}\le1$ is the required final forced price. Its projected
principal polynomial differs from the next prescribed connected
polynomial by at most $y$ on every required healthy region. Purchase
\begin{gather}
 2^n\bar\omega_nR_e y\le d_0,\qquad
 G_i+n H^{-2}R_{e+1}\le H^{-1},\label{eq:allarity-step-purchases}\\
 (2B_i+1)^4H^{-1}\le1,\qquad e^{H^{-1}}-1\le1.
 \label{eq:allarity-step-small}
\end{gather}
The forced inequality is \eqref{eq:allarity-realign-forced-input}
with $s=G_i^2$, $g=g_i$ and $\delta=H^{-1}$. Its output is at most
$G_i$. The regular price becomes at most
\[
       (\alpha+1)z+d_0\le(\alpha+2)z.
\]
Induction therefore bounds the remainder at level $i$ by $d_{k-i}$,
with forced price $G_i$ and the exact prescribed principal polynomial.
At the induction step the density is the push of the preceding marked
density along the same transition kernel; the realignment changes only
its representation. Hence this is an estimate on the original launched
observable, not a sequence of separately chosen approximants.

The number of steps satisfies $(\alpha+2)^{k-l}\le H^{\kappa}$ for a
constant $\kappa$, by Lemma~\ref{lem:allarity-upper-count-power} below.
It follows that $d_{k-l}\le H^{\kappa-q}$.

\subsubsection*{The marker weights}
Two marker weights are used, each the same at every mark, replica and
finite period. Put $H=\mathcal M_k$, where
$\mathcal M_k=\exp\{(\log\log\widehat\beta_k)^2\}$ is the clock monitor
at the comparison depth, as in Section~\ref{sec:rank-root-normalization}.
From the launch at the middle level $\ell$ down to the output torus of
the replica estimate, the weight is
\begin{equation}\label{eq:allarity-marker-weight}
 \omega=H^{1/(8u)} .
\end{equation}
On the leading part of the interval, from the top $k$ down to $\ell$,
the recurrence above is run at the smaller weight
\begin{equation}\label{eq:allarity-upper-weight}
 \omega_-=H^{1/(8u)-U},\qquad U\ge\frac1{8u},
\end{equation}
where the weight exponent $U$ is a constant fixed before the clock scale
(Lemma~\ref{lem:allarity-leading-budget} below says how). Every price
of the recurrence, namely the remainder prices $d_j$, the forced prices
$G_i$, the principal incidence price $H^{-2}$, the discrepancy
contribution \eqref{eq:allarity-realign-regular} and the terminal
prices, is a price in the $\|\cdot\|_{\omega_-}$ norm. The weight is
changed exactly once, at the launch
(Subsection~\ref{subsec:allarity-corrected-budget}), and the cost of
that change is computed there.

Since $U\ge1/(8u)$, we have $\omega_-\le1$, so
\eqref{eq:allarity-mode-weight} gives $\omega_-^{|I|}\le1$ for every
mode. Thus the first purchase in \eqref{eq:allarity-step-purchases} is
implied by
\begin{equation}\label{eq:allarity-weighted-discrepancy}
 2^nR_ey\le d_0=H^{-q}.
\end{equation}
At the large weight, since $H>1$ and $n\le u$,
\begin{equation}\label{eq:allarity-weight-cap}
 \omega^{|I|}\le\bar\omega_n=\omega^n\le\omega^u=H^{1/8}
 \qquad (I\subseteq\{1,\ldots,n\}).
\end{equation}
The incidence and spatial counts, the forced regrouping and the scalar
envelope below hold for every positive weight. The one-step outputs
assumed above, the principal-source bounds and the terminal-source
purchase are hypotheses on prices at the weight $\omega_-$. When such an
input is available as an unweighted bound $b$ on the values of pieces
carrying between one and $n$ markers, with the support charge already
included, it is also a bound on the $\omega_-$-price, because
$\omega_-\le1$.

If the later propagation requires an extra factor $H^{uV+2}$ in the
regular price, purchase $q\ge\kappa+uV+2$, with $\kappa$ the exponent of
Lemma~\ref{lem:allarity-upper-count-power}; this supplies that margin at
the middle level. These two inequalities are separate purchases: a
small initial error does not alone control arbitrarily many steps or an
arbitrary later marker-weight change.

\subsection{Supplying the one-step hypotheses on the leading interval}
\label{subsec:allarity-step-supply}

We now show how the one-step hypotheses of the recurrence follow, on the
leading part of the comparison interval, from the upper calculation of
Section~\ref{sec:actual-mixed-marked-compiler} and from the principal
comparison of Section~\ref{sec:principal-gaussian-comparison}. We also
show that the purchases
\eqref{eq:allarity-step-purchases}--\eqref{eq:allarity-step-small} hold
at large clock scale. Throughout, the top comparison depth is $k$, the
leading interval consists of the rungs $i$ with $\ell\le i<k$, and
$s=\widehat\beta_k$ is the clock scale. We put
\[
 H=\exp((\log\log s)^2)=\mathcal M_k,\qquad
 \omega_-=H^{1/(8u)-U},
\]
as in \eqref{eq:allarity-upper-weight}. The rung couplings $x_i$ (coarse)
and $y_i$ (fine) of rung $i$ lie in $[s/2,s]$. The arity $n\le u$, the
support exponent $e$, the slice multiplier $a$, the support radius
$Q\ge a$, the weight exponent $U$, the tolerance exponent $q\ge6$ and the
step multiplier $\alpha\ge8T_e$ are fixed before $s$, in this order
except that $\alpha$ may be fixed first. Here $T_e$ is the constant of
\eqref{eq:mixed-cover-charge}. A function $f$ of the clock scale has
\emph{upper exponent} $c$ if for every $\varepsilon>0$, eventually
$|f(x)|\le x^{c+\varepsilon}$ uniformly for $x\in[s/2,s]$. In particular
$H^{c}$ has upper exponent $0$ for every real $c$, because
$(\log\log s)^2=o(\log s)$.

\begin{lemma}[Compiler outputs on a leading rung]
\label{lem:allarity-step-outputs}
Assume that, at rung $i$, the upper calculation of
Section~\ref{sec:actual-mixed-marked-compiler} applies at the scale $s$
with the multiplier $\alpha$ and tolerance exponent $q$ fixed above.
Assume the incoming data are:
\begin{itemize}
\item a regular remainder of marker-root price $z$ with
 $H^{-q}\le z\le H^{-4}$;
\item a leading family of marker budget at most $H^{-4}$;
\item a forced background of site price at most $G(y_i)$.
\end{itemize}
All prices are taken at the weight $\omega_-$. Then, after the marked
push, the output has:
\begin{enumerate}
\item nonleading regular marker-root price at most $(\alpha+1)z$;
\item leading marker-root (principal incidence) price at most $H^{-3}\le H^{-2}$;
\item forced site price at most $G(x_i)^2$.
\end{enumerate}
Thus the one-step hypotheses of the recurrence hold at rung $i$ with
$G_i=G(x_i)$.
\end{lemma}
\begin{proof}
In the upper calculation the multipliers are
$l=\max(1,H^{-3}/(\alpha z+H^{-q}))$ and $m=G(x_i)^{-2}$. The incoming
forced budget is $\eta=G(y_i)$, and the leading-family budget is
$h=H^{-4}$. That calculation proves $T_e\eta+a\le H^{-3}$ and
$\Sigma\le K/4$, with $K=H^{-1/8}$. Its rooted connected-set bound then
gives every marker a total marker-root price of at most $T_e\eta+a$ over
the components carrying it. It also gives every site a total site price
of at most $\Sigma\le e^{\Sigma}-1$ over the components containing it.
By \eqref{eq:mixed-club-integral}, the integrated coefficient of a
component is divided by $t_C$:
\begin{itemize}
\item $t_C=1$ for an entirely leading component;
\item $t_C=l$ for an unforced component that is not entirely leading;
\item $t_C=m$ for a component with nonempty forced inventory.
\end{itemize}
Hence:
\begin{enumerate}
\item The nonleading regular components have marker-root price at most
 $H^{-3}/l$. Since $l\ge H^{-3}/(\alpha z+H^{-q})$, this is at most
 $\alpha z+H^{-q}\le\alpha z+z=(\alpha+1)z$, because $z\ge H^{-q}$.
\item The entirely leading components have marker-root price at most
 $T_e\eta+a\le H^{-3}\le H^{-2}$, since $H\ge1$.
\item The forced components have site price at most
 $(e^\Sigma-1)/m\le(e^{K/4}-1)G(x_i)^2\le G(x_i)^2$, because
 $e^{K/4}-1\le K\le1$.
\end{enumerate}
The calculation of that section holds at every positive marker weight,
since its prices include the weights. So these are prices at the weight
$\omega_-$. The multiplier $l$ is at most $\max(1,H^{q-3})$, as in that
section, so the compiler purchase made there covers it.
\end{proof}

The upper bound $z\le H^{-4}$ in the lemma is preserved along the
recurrence. Indeed $d_{k-i}\le d_{k-\ell}\le H^{\kappa-q}\le H^{-4}$ by
Lemma~\ref{lem:allarity-upper-count-power}, once $q\ge\kappa+4$. The
lower bound $z\ge d_0=H^{-q}$ is the convention of the recurrence: a
smaller incoming price may be enlarged to $d_0$.

\begin{lemma}[The leading budget at the small weight]
\label{lem:allarity-leading-budget}
Suppose that at every leading rung the prescribed principal sources have
\emph{unweighted} marker budget at most $H^{A}$, for a constant $A\ge0$
fixed before $s$. That is, for every mark, the sum over the modes
containing it of the value bound times the support charge is at most
$H^{A}$. If
\[
 U\ge A+4+\frac1{8u},
\]
then their marker budget at the weight $\omega_-$ is at most $H^{-4}$.
\end{lemma}
\begin{proof}
Every principal piece has a nonempty marker set $I$, so its weight is
$\omega_-^{|I|}\le\omega_-=H^{1/(8u)-U}\le H^{-A-4}$, because
$\omega_-\le1$. Multiplying the unweighted budget by this factor gives
the claim.
\end{proof}

The launch bounds available for the funded chart bound the principal
source by a fixed positive power of $H$. This is why the leading interval
is run at the weight $\omega_-$: at the large weight $\omega$ the budget
$H^{-4}$ is not available, and at the top comparison level the
connected polynomial $\sum_iz_iP_i$ is not small at all.

\begin{lemma}[A fixed bound for the support charge]
\label{lem:allarity-step-charge}
For a cell $c$ and a nonempty mode $I$, let $r_{c,j}$ be the reader site
of mark $j$ in cell $c$, and define the prescribed support
\[
 X_{c,I}=\{v:\ |v-r_{c,j}|\le QB\ \text{for every } j\in I\}.
\]
Assume that reader sites in distinct replica cells are at distance
greater than $2QB$. Then the geometric hypotheses of the subsection on
marker and spatial prices hold, and every prescribed support satisfies
$\operatorname{cov}_BX_{c,I}\le(\lceil Q\rceil+1)^4$. Hence one may take
\[
 R_e=\exp\{e(\lceil Q\rceil+1)^4\},
\]
a constant fixed before $s$.
\end{lemma}
\begin{proof}
By the first subsection, every point of the full support of a principal
piece of mode $I$ in cell $c$ lies within distance
$(2n+1)\rho\le aB\le QB$ of each of its marked readers, which are the
$r_{c,j}$ with $j\in I$. So the actual support is contained in $X_{c,I}$.
The prescribed support depends only on the mode, by definition. If a site
$v$ lies in $X_{c,I}$ and in $X_{c',I'}$, then a reader of $c$ and a
reader of $c'$ are both within $QB$ of $v$, hence within $2QB$ of each
other. By the separation assumption $c=c'$. Finally, fix one $j\in I$.
Then $X_{c,I}$ is contained in the cube of side $2QB$ centred at
$r_{c,j}$. That cube is covered by a grid of at most
$(\lceil Q\rceil+1)^4$ radius-$B$ cubes whose union is connected. Since
$\operatorname{chg}_e$ is monotone in the cover number, the bound
follows.
\end{proof}

\begin{lemma}[The per-step discrepancy]
\label{lem:allarity-step-discrepancy}
Suppose that at each leading rung the next prescribed connected
polynomial is the Gaussian reference integral of
Section~\ref{sec:principal-gaussian-comparison}, for the same insertion,
forest and tags. Suppose also that the actual projected principal
polynomial is the actual-sector integral there. Suppose finally that,
at each leading rung, the amplitude $\tau$ of the principal insertion has
upper exponent at most $\kappa_\tau<1/2$, uniformly in the rungs. Then,
eventually in $s$, the two polynomials differ by at most $y$ on the
required healthy regions, where $y$ satisfies
\eqref{eq:allarity-weighted-discrepancy}. Hence the first purchase in
\eqref{eq:allarity-step-purchases} holds.
\end{lemma}
\begin{proof}
By \eqref{eq:principal-comparison-final} and the payment paragraph after
it, the difference is at most $\tau$ times a multiplier $\Pi(x)$ of upper
exponent $-1/2$. This holds uniformly in the period and after the finite
forest and tag averages. Take $y=\tau\Pi$, which has upper exponent at
most $\kappa_\tau-1/2<0$. By Lemma~\ref{lem:allarity-step-charge},
$2^nR_e$ is a constant, and $H^{q}$ has upper exponent $0$. Hence
$2^nR_eyH^{q}$ has upper exponent at most $\kappa_\tau-1/2<0$, and is
therefore at most $1$ for all large $s$. This is
\eqref{eq:allarity-weighted-discrepancy}, which implies the first
purchase in \eqref{eq:allarity-step-purchases} because
$\omega_-\le1$.
\end{proof}

\begin{lemma}[The remaining step purchases]
\label{lem:allarity-step-purchases}
For all sufficiently large $s$, uniformly in the rungs of the leading
interval, the following hold:
\begin{enumerate}
\item the second purchase in \eqref{eq:allarity-step-purchases}, with
 $G_i=G(x_i)$;
\item both purchases in \eqref{eq:allarity-step-small}, with
 $B_i=\lceil(\log x_i)^2\rceil^2$.
\end{enumerate}
\end{lemma}
\begin{proof}
\emph{Second purchase.} By Section~\ref{sec:actual-mixed-marked-compiler},
$G$ decays faster than every negative power of $H$ uniformly on
$[s/2,s]$. So eventually $G(x_i)\le H^{-1}/2$. The quantity
$nR_{e+1}$ is a constant by Lemma~\ref{lem:allarity-step-charge}, so
$nR_{e+1}H^{-2}\le H^{-1}/2$ once $H\ge2nR_{e+1}$.

\emph{Smallness.} For $\log s\ge1$ and $x_i\le s$,
$B_i\le((\log s)^2+1)^2\le4(\log s)^4$. Hence
\[
 (2B_i+1)^4\le9^4(\log s)^{16}
 =\exp\{4\log9+16\log\log s\}\le\exp\{(\log\log s)^2\}=H
\]
once $\log\log s\ge17$. Also $e^{H^{-1}}-1\le1$ as soon as
$H^{-1}\le\log2$.
\end{proof}

The step count of the recurrence is
Lemma~\ref{lem:allarity-upper-count-power} below. It uses only $L_j\ge3$
and the maximality of $\ell$; no floor on the block lengths and no stage
condition is needed.

\paragraph{What this supplies and what remains.}
On the leading interval, the lemmas of this subsection and
Lemma~\ref{lem:allarity-upper-count-power} give the following, all at
the weight $\omega_-$ and eventually in $s$ uniformly in the finite
period:
\begin{itemize}
\item the remainder multiplier $(\alpha+1)z$;
\item the principal incidence price $H^{-2}$;
\item the forced price $G_i^2$ with $G_i=G(x_i)$;
\item the leading budget $H^{-4}$, given the unweighted bound $H^A$;
\item the discrepancy purchase \eqref{eq:allarity-weighted-discrepancy};
\item the purchases
 \eqref{eq:allarity-step-purchases}--\eqref{eq:allarity-step-small};
\item the step count, with remainder $d_{k-\ell}\le H^{\kappa-q}$ at the
 middle level.
\end{itemize}
The following inputs are not supplied here.
\begin{enumerate}
\item \emph{Applicability of the upper calculation} at each leading rung.
 This includes its primitive body-compiler, all-hole and joint-warning
 purchases, which Section~\ref{sec:actual-mixed-marked-compiler} lists as
 inputs on the same body alphabet and coarse representation.
\item \emph{The unweighted principal-source bound} $H^A$ of
 Lemma~\ref{lem:allarity-leading-budget} at every leading rung, and the
 terminal-source purchase: the launched insertions must give the initial
 remainder price $d_0$ and forced price $G_k$ at the weight $\omega_-$.
\item \emph{The identification in Lemma~\ref{lem:allarity-step-discrepancy}}
 of the prescribed connected polynomial with the Gaussian reference
 integral, and the bound $\kappa_\tau<1/2$ on the upper exponent of the
 amplitude at every leading rung.
\item \emph{The reader separation} $2QB$ between distinct replica cells
 assumed in Lemma~\ref{lem:allarity-step-charge}.
\end{enumerate}
The rungs below the middle level use the lower calculation of
Section~\ref{sec:actual-mixed-marked-compiler} at the large weight
$\omega$, after the launch of
Subsection~\ref{subsec:allarity-corrected-budget}.

\subsection{The scalar envelope used in the replica estimate}

We record the final elementary step linking a paid marked representation
to Section~\ref{sec:rank-root-normalization}. Suppose a polynomial
density on a torus of $q_0$ sites has the form
\[
 D(U)=e^{\operatorname{exponent}_P(U)}
       \exp_*\left(\sum_a z_{I_a}h_a(U)\right)
       \sum_{B\ {\rm admissible}}\prod_{b\in B}K_b(U).
\]
The $I_a$ are nonempty, and each background packet has nonempty full
spatial support. Let $\|\cdot\|_{\boldsymbol r}$ denote the coefficient
$\ell^1$ norm weighted by positive marker weights. Assume regular
incidence price at most $\ell$ per mark and forced spatial price at
most $L$ per site, with nonnegative support exponent. If there are
$M_0$ marks in total, double counting gives
\[
 \sum_a\|z_{I_a}h_a\|_{\boldsymbol r}\le M_0\ell,
       \qquad\sum_b\|K_b\|_{\boldsymbol r}\le q_0 L.
\]
A nonempty support contributes at least once to the corresponding
incidence sum. Dropping the support weight is legitimate because it
is at least one. Masks can only lower the regular value bounds.
Submultiplicativity gives
\[
 \left\|\exp_*\left(\sum_a z_{I_a}h_a\right)\right\|_{\boldsymbol r}
       \le e^{M_0\ell}.
\]
Dropping the admissibility restrictions in the nonnegative majorant
bounds the forced sum by
$\prod_b(1+\|K_b\|_{\boldsymbol r})\le e^{q_0 L}$. Therefore
\begin{equation}\label{eq:allarity-density-envelope}
 \|D(U)\|_{\boldsymbol r}
      \le\exp\{\operatorname{exponent}_P(U)+q_0L+M_0\ell\}.
\end{equation}
For the replica estimate, first complete any later propagation to the
selected output torus whose site count is exactly the number $q_0$ of
replica cells. This equality is a property of that selected output
level, not of an arbitrary intermediate torus. At this output there
are $M_0=q_0 n$ marks. Thus $L,\ell\le1$
give per-site expense at most $1+n$, exactly the arity dependence
in the determinant ticket. The weight of the coefficient carrying
a chosen nonempty subset $I$ in each of the $q_0$ cells is the
$q_0$th power of its single-cell weight. At the constant weight
\eqref{eq:allarity-marker-weight} this is $\omega^{q_0|I|}$, and
since the coefficient norm dominates each weighted coefficient,
\eqref{eq:allarity-density-envelope} gives pointwise, for the
coefficient $D_{I}^{\otimes}(U)$ of the monomial carrying $I$ in every
cell,
\[
 \omega^{q_0|I|}\,\bigl|D_{I}^{\otimes}(U)\bigr|
   \le\exp\{\operatorname{exponent}_P(U)+q_0(L+n\ell)\}.
\]
With $u=M$ the marker allowance of
Section~\ref{sec:rank-root-normalization}, $\omega=\mathcal
M_k^{1/(8M)}=e^{2I_k}$ is the weight $a$ used there, because
$I_k=\log\mathcal M_k/(16M)$. Hence this is the coefficient price
$a^{q_0|I|}$ with per-site expense $X=L+n\ell\le1+M$ required there,
on the same pushed density. Apply the normalization
estimate and then take the $q_0$th root. The denominator is the
partition function of the same pushed density, preserved by the
unmarked coefficient identity at every step.

The recurrence above states its one-step hypotheses quantitatively.
The joint Gaussian reader estimate supplies the discrepancy $y$ and
the principal source bounds; the mixed compiler supplies the exact
marked push and the pre-realignment prices. The selected scale
construction must buy all of these on the common history, including
the later propagation below $l$. Neither the recurrence nor the
scalar envelope asserts those upstream estimates without their
corresponding purchases.

\subsection{Transport to the output torus}
\label{subsec:allarity-output-transport}

We complete the propagation in the application to the pairwise receipt.
The chart has level $k$, middle level $\ell$ and root depth $w\le\ell$, and
the arity is $n=2$. We identify the marked density at the output, show
that its pushed version below the middle level keeps prices at most one,
and conclude with the coefficient price of
Section~\ref{sec:rank-root-normalization}. As before, $s=\widehat\beta_k$,
$H=\mathcal M_k$, $\omega=H^{1/(8u)}$, and the fixed constants
$\alpha\ge8T_e$ and $q\ge6$ are those of
Subsection~\ref{subsec:allarity-step-supply}.

\paragraph{The invoice and its parts.}
Let $f_0,f_1$ be the two finite-volume insertions of a fit configuration,
as in Subsection~\ref{subsec:trace-rounded-receipt}. Let $S_0(T)$
($T\subseteq\{0,1\}$) be the reference moments at the rounded
configuration, with $S_0(\varnothing)=1$. In the marker algebra
$\mathscr A_{\{0,1\}}$ of Section~\ref{sec:allarity-planted-compiler},
the polynomial $\sum_TS_0(T)z_T$ has constant coefficient one. So it is
$\exp_*D$ for a unique $D$ with zero constant coefficient, namely
\[
 D=j_0z_0+j_1z_1+\kappa z_0z_1,\qquad
 j_c=S_0(\{c\}),\quad\kappa=S_0(\{0,1\})-j_0j_1 .
\]
Indeed $\exp_*D=1+D+D^2/2$ and $D^2=2j_0j_1z_0z_1$. The \emph{invoice} is
\[
 \mathcal I=\exp_*(-D)\,(1+f_0z_0)(1+f_1z_1).
\]

\begin{lemma}[The parts of the invoice]\label{lem:allarity-invoice-parts}
The coefficients of $\mathcal I$ are
\[
 \mathcal I_{\{c\}}=f_c-j_c\quad(c=0,1),\qquad
 \mathcal I_{\{0,1\}}=(f_0-j_0)(f_1-j_1)-\kappa ,
\]
that is, $\mathcal I_s=\Phi_s$ for every nonempty $s$.
\end{lemma}
\begin{proof}
First,
$\exp_*(-D)=1-D+D^2/2=1-j_0z_0-j_1z_1+(j_0j_1-\kappa)z_0z_1$.
Multiply by $1+f_0z_0+f_1z_1+f_0f_1z_0z_1$, using $z_c^2=0$. The
coefficient of $z_c$ is $f_c-j_c$. The coefficient of $z_0z_1$ is
\[
 f_0f_1-j_0f_1-j_1f_0+j_0j_1-\kappa
 =(f_0-j_0)(f_1-j_1)-\kappa .\qedhere
\]
\end{proof}

\paragraph{Cells and sites.}
The cells of the centered tiling at depth $w$ of the level-$N$ torus are
indexed by the sites of the level-$w$ torus, one cell for each site.
Hence the number $q_0$ of replica cells equals the site count of the
level-$w$ torus. This is the equality required in the scalar envelope:
the output torus is the level-$w$ torus.

\paragraph{Propagation below the middle level.}
On the rungs $j$ with $w\le j<\ell$ there are no leading letters, and the
lower calculation of Section~\ref{sec:actual-mixed-marked-compiler}
applies. All prices from the launch at $\ell$ downwards are taken at the
large weight $\omega=H^{1/(8u)}$. With incoming regular marker-root price
$z\le H^{-1/2}$ and incoming forced budget $z\,\omega_j$, where
$\omega_j\in[0,1]$ is the actual lower-interval attenuation, the lower
calculation gives two outputs:
\begin{itemize}
\item regular marker-root price at most
 $\tfrac12(\alpha z+H^{1/2}z^2)\le\alpha z+H^{1/2}z^2$;
\item forced site price at most $(e^{\Sigma}-1)/m\le G(y_j)$, where
 $m=G(y_j)^{-1}$ and $y_j$ is the fine coupling of the rung.
\end{itemize}
These require the rung couplings to lie in $[s/2,s]$.

\begin{lemma}[Lower propagation]\label{lem:allarity-lower-propagation}
Assume the following.
\begin{enumerate}
\item At level $\ell$ the regular marker-root price is at most
 $z_0\le H^{-3/4}$, and there are no leading letters.
\item Every rung of $[w,\ell)$ satisfies the hypotheses of the lower
 calculation, with couplings in $[s/2,s]$.
\item The forced input of each lower rung is admissible:
 its forced budget is at most $z_t\,\omega_j$ along the iteration below.
 Prices are upper bounds, so $z_t$ may be enlarged for this.
\item The lower step count satisfies
 \begin{equation}\label{eq:allarity-lower-count}
  (\alpha+1)^{\ell-w}\le H^{1/4}.
 \end{equation}
\end{enumerate}
Let $z_t$ be the regular bound after $t$ lower steps, so that level
$\ell-t$ carries $z_t$. Then $z_t\le(\alpha+1)^tz_0\le H^{-1/2}\le1$ for
$0\le t\le\ell-w$, and every forced site price on $[w,\ell]$ is at most
$1$.
\end{lemma}
\begin{proof}
If $H^{1/2}z_t\le1$, then
\[
 z_{t+1}\le(\alpha+H^{1/2}z_t)z_t\le(\alpha+1)z_t .
\]
By induction $z_t\le(\alpha+1)^tz_0$ for $t\le\ell-w$, as long as
$H^{1/2}z_t\le1$. This condition holds throughout, because
$(\alpha+1)^tz_0\le H^{1/4}H^{-3/4}=H^{-1/2}$ by
\eqref{eq:allarity-lower-count}. That is also the domain bound
$z\le H^{-1/2}$ of the lower calculation. The forced outputs of the lower
rungs are at most $G(y_j)\le1$. At $t=0$ the forced price is the one
delivered by the launch, which is at most $1$ by
Proposition~\ref{prop:allarity-launch}.
\end{proof}

By Lemma~\ref{lem:allarity-lower-count-free} below, the lower count
\eqref{eq:allarity-lower-count} holds for all large $s$, uniformly in the
chart, and by Lemma~\ref{lem:allarity-lower-forced} hypothesis~3 can be
dropped.

\begin{proposition}[The coefficient price at the output]
\label{prop:allarity-output-price}
Assume the following.
\begin{enumerate}
\item The launch receipt of Proposition~\ref{prop:allarity-launch} at the
 middle level.
\item The hypotheses of Lemma~\ref{lem:allarity-lower-propagation}.
\item The representation at the output level $w$ is lawful, with the
 suppositions of Section~\ref{sec:rank-root-normalization} at the stage
 of level $w$.
\end{enumerate}
Then, at the weight $\omega=H^{1/(8u)}$, the following hold for every
nonempty $s\subseteq\{0,1\}$:
\[
 \omega^{q_0|s|}\,\bigl|\mathcal D_s^{\otimes}(U)\bigr|
 \le\exp\{\operatorname{exponent}_P(U)+q_0(1+n)\}.
\]
Here $\mathcal D$ is the push to the level-$w$ torus of the replicated
marked density of the invoice $\mathcal I$, and $\mathcal D_s^{\otimes}$
is its coefficient carrying $s$ in every cell. The normalizing
denominator is the partition function of the same pushed density, by the
unmarked coefficient identity at every step. Taking $u=M$, so that
$\omega=a$, this is the coefficient price $a^{q_0|s|}$ with per-site
expense $X\le1+n\le1+M$ required by \eqref{eq:rank-root}. Consequently,
for $u=M$ and $k\ge K_M$ of Proposition~\ref{prop:rank-root-output-ticket},
the root test of $\Phi_s$ on the centered tiling at depth $w$ is at most
$\mathcal M_k^{-1/(16M)}$.
\end{proposition}
\begin{proof}
By Lemma~\ref{lem:allarity-lower-propagation}, at level $w$ the regular
incidence price per mark is $\ell_w\le1$ and the forced site price is
$L_w\le1$. By the cells-and-sites paragraph, the output torus has $q_0$
sites, and there are $M_0=q_0n$ marks. The scalar envelope
\eqref{eq:allarity-density-envelope} and the paragraph after it give the
displayed inequality with expense $L_w+n\ell_w\le1+n$. By
Lemma~\ref{lem:allarity-invoice-parts}, the coefficients of the
invoice are the $\Phi_s$, so the replicated coefficient is that of
$\Phi_s$ in every cell. The last statement is
Proposition~\ref{prop:rank-root-output-ticket}, applied with $w(k)=w$,
together with \eqref{eq:rank-root} and the lawful representation at $w$.
For general $u$ the weight is $\omega=H^{1/(8u)}$ and the same argument
gives the bound $H^{1/(16M)-|s|/(8u)}$ only; the root conclusion uses
$u=M$.
\end{proof}

\subsection{The launch at the middle level}
\label{subsec:allarity-corrected-budget}

At the middle level $\ell$ the representation produced by the leading
interval is converted into one suitable for the lower interval. Three
things change at once: the marker weight is raised from $\omega_-$ to
$\omega$; the leading family is replaced by its corrected form; and the
leading designation is dropped, so that the corrected pieces become
ordinary regular letters. We keep the notation of
Subsections~\ref{subsec:allarity-step-supply}
and~\ref{subsec:allarity-output-transport}. In particular
$\mathcal I=\exp_*(-D_R)\prod_c(1+f_cz_c)$ is the invoice, with $D_R=D$ the
reference cumulant polynomial.

\paragraph{The reference factor commutes with every push.}
The polynomial $\exp_*(-D_R)$ has field-independent coefficients. Every
block integration and every push acts linearly and coefficientwise on
the marker algebra. So each of them commutes with multiplication by
$\exp_*(-D_R)$. Hence the pushed invoice at any level is
$\exp_*(-D_R)$ times the pushed product of the insertions. At the middle
level the latter has principal part $\exp_*(D_Q)$, with $D_Q$ the
connected polynomial of the same Gaussian history. The invoice is
therefore launched at the middle level with corrected principal
polynomial $D_Q-D_R$. This is an identity of densities. It does not by
itself give a priced representation, for the following reason.

\paragraph{The failed mask.}
For one mode, let $\chi$ be the healthy mask of its prescribed support,
$W$ the actual leading coefficient and $R$ the corresponding coefficient
of $D_R$. The leading letter carries $\chi W$, while the reference factor
contributes $R$ everywhere. Hence
\[
 \chi W-R=\chi(W-R)-(1-\chi)R .
\]
The first term is a regular letter controlled by the healthy bound on
$W-R$ (Lemma~\ref{lem:allarity-corrected-pieces}). The second term is
supported where the prescribed support contains a flag. It must be moved
into the forced alphabet and regrouped with the forced background, as in
the mask change of the realignment subsection. Its price is controlled by
a bound on the reference coefficients $R$, not by the healthy bound.

\paragraph{The corrected budget.}
The bound on the corrected pieces is already proved in
Section~\ref{sec:allarity-cumulant-comparison}. There the free moment
polynomial $q'$ is the reference moment polynomial $1+X_R$, and the
moment polynomial $q$ of the same Gaussian history is $1+X_Q$. If the
coefficientwise moment error is at most $H^{-3/2}$, then
\eqref{eq:allarity-cumulant-error} gives
$\|D_Q-D_R\|_1+H^{-3/2}\le H^{-31/32}$ for all large $s$. That moment
error is \eqref{eq:allarity-moment-error} at the powers of its scale
paragraph. The actual healthy local error between a planted leading
coefficient and the coefficient of $D_Q$ is at most $H^{-3/2}$ by
Section~\ref{sec:marked-healthy-cumulant-launch}. Both statements hold
under the hypotheses stated in those sections.

\begin{lemma}[Corrected leading pieces]
\label{lem:allarity-corrected-pieces}
Assume the moment error $H^{-3/2}$ of
Section~\ref{sec:allarity-cumulant-comparison} and the healthy local error
$H^{-3/2}$ of Section~\ref{sec:marked-healthy-cumulant-launch}, at the
middle level. Then, for all large $s$, every corrected leading piece at
that level, that is, an actual leading piece minus the corresponding
coefficient of $D_R$, is at most $H^{-31/32}$ on its healthy region. Its
incidence price at the weight $\omega$ is at most
\[
 h=c\,H^{-27/32},\qquad c=2^nR_e .
\]
\end{lemma}
\begin{proof}
The first assertion is the triangle inequality applied to the healthy
local error and \eqref{eq:allarity-cumulant-error}, as in the last
paragraph of Section~\ref{sec:allarity-cumulant-comparison}. Each
coefficient is bounded by the coefficient $\ell^1$ norm. For the price, a
piece of mode $I$ carries weight $\omega^{|I|}\le H^{1/8}$ by
\eqref{eq:allarity-weight-cap}, and support charge at most $R_e$ by
Lemma~\ref{lem:allarity-step-charge}. At most $2^n$ modes contain a given
mark. Hence the incidence price is at most
$2^nR_eH^{1/8}H^{-31/32}=cH^{-27/32}$.
\end{proof}

\begin{lemma}[Raising the weight]
\label{lem:allarity-weight-raise}
Let a family of regular letters, each carrying a marker set $I$ with
$|I|\le n$, have marker-root or incidence price $p$ at the weight
$\omega_-$. Then its price of the same kind at the weight $\omega$ is at
most $H^{Un}p$.
\end{lemma}
\begin{proof}
These prices are sums of nonnegative terms, one for each letter, each
proportional to the weight $\omega_-^{|I|}$ of that letter. Raising the
weight multiplies the term by $(\omega/\omega_-)^{|I|}=H^{U|I|}\le
H^{Un}$, since $H>1$.
\end{proof}

The forced background is not covered by this lemma: a forced component
may contain marked pieces from several cells, and its weight may then
grow by more than $H^{Un}$. Its weight change is part of the following
input.

The launch requires one further input, which this section does not prove.
\begin{enumerate}
\item[(L)] \emph{Failed-mask regrouping and weight change at the launch.}
 There is a constant $c_L\ge0$, fixed before $s$, with the following
 property. At the middle level, let the forced background have site price
 at most $\eta_-$ at the weight $\omega_-$. Move the failed-mask pieces
 $-(1-\chi)R\,z_I$ of all modes into the forced alphabet and regroup them
 with that background. Then the result is an exact representation of the
 same density, whose regular letters are unchanged and whose forced site
 price at the weight $\omega$ is at most $H^{c_L}\eta_-$.
\end{enumerate}
The input (L) needs a bound on the reference coefficients $R$ by a fixed
power of $H$, control of the forced background at the larger weight, and
the forced regrouping of Section~\ref{sec:soft-log-fusion} with an
inflation exponent lowered by a multiple of $\log H$.

\begin{proposition}[The launch receipt]
\label{prop:allarity-launch}
Assume the conclusions of Subsection~\ref{subsec:allarity-step-supply} on
the leading interval, with remainder $d_{k-\ell}\le H^{\kappa-q}$ and
forced price $G(x_\ell)$ at the weight $\omega_-$. Assume
Lemma~\ref{lem:allarity-corrected-pieces} at the middle level and the
input (L). Assume the tolerance exponent satisfies
\[
 q\ge\kappa+Un+1 .
\]
Then, for all large $s$, the pushed invoice at the middle level has an
exact representation at the weight $\omega$, without leading letters, with
regular marker-root price at most $H^{-3/4}$ and forced site price at most
$H^{-q}$.
\end{proposition}
\begin{proof}
By the commuting reference factor, the pushed invoice at $\ell$ is the
pushed insertion product times $\exp_*(-D_R)$. The representation from
the leading interval has three parts. Their prices are as follows.

\emph{Remainder.} By Lemma~\ref{lem:allarity-weight-raise}, the regular
remainder has price at most $H^{Un}H^{\kappa-q}\le H^{-1}$ at the weight
$\omega$.

\emph{Leading family.} Its pieces are the exact prescribed principal
polynomial. Multiplying by $\exp_*(-D_R)$ and splitting each corrected
coefficient by the failed-mask identity gives healthy pieces
$\chi(W-R)$ and failed-mask pieces. Drop the leading designation. The
healthy pieces become regular letters of incidence price at most
$cH^{-27/32}$, by Lemma~\ref{lem:allarity-corrected-pieces}.

\emph{Forced part.} The forced background has site price at most
$\eta_-=G(x_\ell)$ at the weight $\omega_-$. This is the realigned forced
price of the recurrence, not the compiler output $G(x_\ell)^2$ before
realignment. By (L), after the failed-mask pieces are moved and
regrouped, the forced site price at the weight $\omega$ is at most
$H^{c_L}G(x_\ell)$.

Adding the regular contributions gives at most $H^{-1}+cH^{-27/32}$, which
is at most $H^{-3/4}$ once $H^{3/32}\ge c+1$. Since $G$ decays faster than
every negative power of $H$ uniformly on $[s/2,s]$, eventually
$H^{c_L}G(x_\ell)\le H^{-q}$.
\end{proof}

\subsection{The step counts without a block-length floor}
\label{subsec:allarity-lower-count}

\begin{lemma}[The lower step count is free]
\label{lem:allarity-lower-count-free}
Let the clock's cover size at level $k$ be
$c_k=\lceil(\log\widehat\beta_k)^2\rceil$, and let the root depth be
$w=\max\{j\le\ell:a_j/a_\ell\ge c_k^{10}\}$. Since all block lengths
satisfy $L_j\ge3$, for every fixed $\gamma>0$ and all sufficiently large
$s=\widehat\beta_k$, uniformly in the chart,
$(\alpha+1)^{\ell-w}\le H^{\gamma}$. In particular the lower step count
\eqref{eq:allarity-lower-count} holds.
\end{lemma}
\begin{proof}
If $w=\ell$ there is nothing to prove. Otherwise, by maximality of $w$,
$a_{w+1}/a_\ell<c_k^{10}$. By Section~\ref{sec:common-scale-selection},
$a_{w+1}/a_\ell=\prod_{t=w+2}^{\ell}L_{J(t)}\ge3^{\ell-w-1}$. For
$\log s\ge1$ we have $c_k\le2(\log s)^2$. Hence
\[
 (\ell-w-1)\log3<10\log c_k\le10\log2+20\log\log s .
\]
Therefore
\[
 (\ell-w)\log(\alpha+1)\le\log(\alpha+1)
   \Bigl(1+\frac{10\log2+20\log\log s}{\log3}\Bigr)
 =A_0+A_1\log\log s ,
\]
with constants $A_0,A_1\ge0$ fixed by $\alpha$. Since $\gamma>0$ and
$\log H=(\log\log s)^2$, the right side is at most $\gamma\log H$ for all
large $s$. No block-length floor and no stage condition is used.
\end{proof}

\begin{lemma}[The upper step count at a fixed power]
\label{lem:allarity-upper-count-power}
Let the base arc satisfy $b\le\bar b$, and put
\[
 \kappa=1+\frac{(\bar b+2)\log(\alpha+2)}{\log3}.
\]
For all sufficiently large $s$, uniformly in the chart,
$(\alpha+2)^{k-\ell}\le H^{\kappa}$. Consequently the remainder price of
the recurrence at the middle level is $d_{k-\ell}\le H^{\kappa-q}$ at the
weight $\omega_-$. If $q\ge\kappa+4$, this is at most $H^{-4}$, which is
the window required in Lemma~\ref{lem:allarity-step-outputs}.
\end{lemma}
\begin{proof}
If $\ell=k$ there is nothing to prove. Otherwise
$a_{\ell+1}/a_k<m\mu_k^2=H^{b+2}$ by maximality of $\ell$. Also
$a_{\ell+1}/a_k=\prod_{t=\ell+2}^{k}L_{J(t)}\ge3^{k-\ell-1}$. Hence
$(k-\ell-1)\log3<(b+2)\log H$, and
\[
 (k-\ell)\log(\alpha+2)\le\log(\alpha+2)
   +\frac{(\bar b+2)\log(\alpha+2)}{\log3}\log H\le\kappa\log H
\]
once $\log H\ge\log(\alpha+2)$. Then
$d_{k-\ell}=(\alpha+2)^{k-\ell}H^{-q}\le H^{\kappa-q}$.
\end{proof}

The parameters are therefore fixed in the order $\alpha$, $\bar b$, then
$U$ (by Lemma~\ref{lem:allarity-leading-budget}), then
\[
 q\ge\max\{6,\ \kappa+4,\ \kappa+Un+1\},
\]
which covers Lemma~\ref{lem:allarity-step-outputs} and
Proposition~\ref{prop:allarity-launch}. All are fixed before $s$.

\begin{lemma}[The forced inputs of the lower rungs]
\label{lem:allarity-lower-forced}
In Lemma~\ref{lem:allarity-lower-propagation}, hypothesis~3 can be
dropped. More precisely, take the incoming forced budget of each lower
rung to be the forced output $G(y)$ of the preceding lower rung, and at
the transition the forced output of the launch, which is at most $H^{-q}$
by Proposition~\ref{prop:allarity-launch}. Replace the bounds $z_t$ by
$\max(z_t,H^{-q})$. Then the conclusions of that lemma hold for all large
$s$.
\end{lemma}
\begin{proof}
The lower calculation of Section~\ref{sec:actual-mixed-marked-compiler}
uses its incoming forced budget $\eta$ in exactly two places: the
marker-root term $T_e\eta$, and the site term
$\mathcal W(B_c)L^4T_e\eta$. With $\eta=z\omega_j\le z$, it bounds them by
$T_ez$ and by $K/16$ for $z\le H^{-1/2}$. Both bounds hold verbatim for
any $\eta\le z$. Now $G$ decays faster than every power of $H$ uniformly
on $[s/2,s]$, so for large $s$ every incoming forced budget of a lower
rung is at most $H^{-q}\le\max(z_t,H^{-q})$; at the transition this is
the launch bound. Replacing $z_t$ by $\max(z_t,H^{-q})$ keeps the
induction of Lemma~\ref{lem:allarity-lower-propagation}: if
$H^{-q}\le z_t\le H^{-1/2}$, then
$\max(\alpha z_t+H^{1/2}z_t^2,H^{-q})\le(\alpha+1)z_t$, since
$\alpha\ge1$. The starting value is $\max(z_0,H^{-q})\le H^{-3/4}$. The
forced outputs are at most $G(y_j)\le1$ as before.
\end{proof}

\begin{lemma}[The amplitude at the middle level]
\label{lem:allarity-amplitude}
Under the hypotheses of Lemma~\ref{lem:allarity-corrected-pieces} there is
a constant $C_*$, depending on the fixed word and the compact window,
such that for all large $s$ every actual leading piece at the middle level
is bounded by $C_*$ on its healthy region.
\end{lemma}
\begin{proof}
The reference moment polynomial $1+X_R$ varies continuously on the
compact window, which consists of configurations with distinct points.
So $\|X_R\|_1\le M_{\mathcal W}$ for a constant $M_{\mathcal W}$, and
$\|D_R\|_1\le\sum_{j=1}^n M_{\mathcal W}^j/j$ by submultiplicativity. By
Lemma~\ref{lem:allarity-corrected-pieces}, each actual leading piece is
within $H^{-31/32}\le1$ of a coefficient of $D_R$. Take
$C_*=1+\sum_{j=1}^nM_{\mathcal W}^j/j$.
\end{proof}

The lemma bounds the amplitude at the middle level only. The discrepancy
lemma needs the amplitude bound $\kappa_\tau<1/2$ at every leading rung,
which is listed below as an input. Any bound $CH^{A'}$ with fixed $A'$
would suffice there, since such a bound has upper exponent $0$.

\paragraph{Status of the propagation to the output.}
Proposition~\ref{prop:allarity-output-price} is proved from the
following inputs, all on the common history of the chart:
\begin{enumerate}
\item the four inputs listed at the end of
 Subsection~\ref{subsec:allarity-step-supply} for the leading interval:
 applicability of the upper calculation; the unweighted principal-source
 bound $H^A$ and the terminal-source purchase; the identification of the
 prescribed polynomial and the per-rung amplitude $\kappa_\tau<1/2$; the
 reader separation $2QB$;
\item the two $H^{-3/2}$ inputs of Lemma~\ref{lem:allarity-corrected-pieces}
 at the middle level, which are proved in
 Sections~\ref{sec:allarity-cumulant-comparison}
 and~\ref{sec:marked-healthy-cumulant-launch} under their stated
 hypotheses (the reader covariance and vertex bounds, the exact low-grade
 calibration of Section~\ref{sec:allarity-coordinate-calibration}, the
 mixed compiler's per-step row and covariance comparison, the terminal
 paid insertion bound and the common purchases);
\item the failed-mask regrouping (L) at the launch;
\item applicability of the lower calculation on $[w,\ell)$;
\item the lawful representation at the output level $w$, with the
 suppositions of Section~\ref{sec:rank-root-normalization}.
\end{enumerate}
The rung couplings lie in $[s/2,s]$ on the whole window $[w,k]$ by the
plateau field of the chart witness, which is proved from the clock in
Lemma~\ref{lem:ex-lattice-firesky-clock-fields}; the same lemma proves the
bracket. The step counts, the forced inputs of the lower rungs and the
leading budget at the small weight are proved above. The polynomial
transport purchase of Section~\ref{sec:marked-healthy-cumulant-launch} is
separate from the step counts and remains one of the common purchases.
